\documentclass[12pt,a4paper]{article}

% =========================================================
% QGX Dissertation Template (Rule-Compliant)
% =========================================================

\usepackage[utf8]{inputenc}
\usepackage[T1]{fontenc}
\usepackage{geometry}
\usepackage{setspace}
\usepackage{fancyhdr}
\usepackage{titlesec}
\usepackage{indentfirst}
\usepackage{graphicx}
\usepackage{booktabs}
\usepackage{tabularx}
\usepackage{amsmath}
\usepackage{amssymb}
\usepackage{hyperref}
\usepackage{tocbibind}
\usepackage{caption}
\usepackage{float}
\usepackage{times}

% Margins as requested
\geometry{
  a4paper,
  left=1.5in,
  right=1in,
  top=1in,
  bottom=1in,
  gutter=0in
}

% Hyperlinks
\hypersetup{colorlinks=true, linkcolor=black, citecolor=black, urlcolor=blue}

% Paragraph and line spacing rules
\onehalfspacing
\setlength{\parindent}{0.9cm}
\setlength{\parskip}{0pt}

% Caption style
\captionsetup{font={bf,normalsize}, justification=centering}

% Section heading formats
\titleformat{\section}
  {\bfseries\fontsize{16}{18}\selectfont}
  {Chapter \thesection}
  {0.8em}
  {}

\titleformat{\subsection}
  {\bfseries\fontsize{14}{16}\selectfont}
  {\thesection.\arabic{subsection}}
  {0.8em}
  {}

\titleformat{\subsubsection}
  {\bfseries\fontsize{12}{14}\selectfont}
  {\thesection.\arabic{subsection}.\arabic{subsubsection}}
  {0.8em}
  {}

% Header/footer configuration
\pagestyle{fancy}
\fancyhf{}
\fancyhead[L]{\nouppercase{\leftmark}}
\fancyfoot[R]{\thepage}
\renewcommand{\headrulewidth}{0.4pt}
\renewcommand{\footrulewidth}{0pt}
\renewcommand{\sectionmark}[1]{\markboth{Chapter \thesection: #1}{}}

% Utility command for placeholders
\newcommand{\placeholderfigure}[2]{%
\begin{figure}[H]
\centering
\fbox{\parbox[c][2.2in][c]{0.85\textwidth}{\centering\textbf{#2}}}
\caption{#2}
\label{fig:#1}
\end{figure}
}

\begin{document}

% =========================================================
% FRONT MATTER (Roman page numbers, bottom-middle)
% =========================================================

\pagenumbering{roman}
\fancyhf{}
\fancyfoot[C]{\thepage}
\fancyhead{}
\renewcommand{\headrulewidth}{0pt}

\begin{titlepage}
\centering
\vspace*{1.5cm}
{\bfseries\fontsize{18}{22}\selectfont QGX-D1\par}
\vspace{0.7cm}
{\bfseries\fontsize{16}{18}\selectfont A Cloud-Native AI-Enabled Learning Management System\par}
\vspace{1.2cm}
{\bfseries\fontsize{14}{16}\selectfont PROJECT DISSERTATION\par}
\vspace{1.2cm}
Submitted in partial fulfillment of the requirements for the award of degree\\
\vspace{0.7cm}
\textbf{Bachelor of Engineering / Technology}\\
\vspace{1.5cm}
\textbf{Candidate Name:} ............................................................\\
\vspace{0.4cm}
\textbf{Register Number:} ............................................................\\
\vspace{0.4cm}
\textbf{Department:} ............................................................\\
\vfill
\textbf{Institution Name}\\
\textbf{Month Year}
\end{titlepage}

\newpage
\section*{Certificate}
This is to certify that the dissertation entitled \textbf{QGX-D1: A Cloud-Native AI-Enabled Learning Management System} is a bona fide record of work carried out by the candidate under our supervision and guidance.

\vspace{2cm}
\noindent\textbf{Project Guide Signature:} \hrulefill
\vspace{1.2cm}
\noindent\textbf{Head of Department Signature:} \hrulefill
\vspace{1.2cm}
\noindent\textbf{Principal Signature:} \hrulefill

\newpage
\tableofcontents
\newpage
\listoffigures
\newpage
\listoftables
\newpage

% =========================================================
% MAIN MATTER (Arabic page numbers, bottom-right)
% =========================================================

\pagenumbering{arabic}
\fancyhf{}
\fancyhead[L]{\nouppercase{\leftmark}}
\fancyfoot[R]{\thepage}
\renewcommand{\headrulewidth}{0.4pt}

\section{Introduction}

\subsection{Purpose of the System}
QGX-D1 is designed to provide an integrated learning management platform for institutional academic operations. The purpose of the system is to unify authentication, role-based operations, AI-assisted query support, monitoring, and communication workflows in one secure application. The system improves administrative efficiency, data consistency, and academic process transparency.

\subsection{Scope of the System}
The scope includes user onboarding, secure login, role-specific dashboards, assessment support, notification workflows, audit logs, password recovery, and offline-ready progressive web capabilities. The system targets four primary actors: Admin, Teacher, Student, and Parent. The deployment scope includes cloud-native hosting and managed database infrastructure.

\subsection{Existing System and Limitations}
Existing institutional workflows are often fragmented across spreadsheets, messaging apps, independent portals, and manual records. This produces duplicate data, delayed communication, weak access control consistency, and limited auditability. Similar issues have also been observed in educational platform studies \cite{anderson2022,lee2023}.

\subsection{Proposed System and Advantages}
QGX-D1 proposes a unified architecture with centralized identity, role governance, workflow tracking, and AI support. Advantages include reduced manual work, stronger role enforcement, real-time visibility, consistent user experience, and improved readiness for institutional scaling.

\placeholderfigure{intro-context}{Figure 1.1: QGX-D1 System Context Overview}

\section{Requirement Analysis}

\subsection{Functional Requirements}
Key functional requirements are listed below.
\begin{enumerate}
\item User registration, login, logout, and password recovery.
\item Role assignment and role-based navigation.
\item Dashboard operations for Admin, Teacher, Student, and Parent.
\item AI query generation with controlled backend processing.
\item Batch user creation for administrative onboarding.
\item Activity log generation for critical system actions.
\item Notification delivery and user feedback handling.
\item PWA installation and offline fallback support.
\end{enumerate}

\subsection{Non-Functional Requirements}
\begin{enumerate}
\item Security with authentication, authorization, validation, and logging.
\item Performance with responsive page loads and efficient API latency.
\item Reliability with graceful failure handling and recovery paths.
\item Usability with role-aware and accessible interface design.
\item Maintainability with modular code and documented architecture.
\item Scalability for institutional growth and concurrent users.
\end{enumerate}

\subsection{Hardware Requirements}
\begin{itemize}
\item Processor: Dual-core or above.
\item RAM: Minimum 4 GB (8 GB recommended).
\item Storage: 10 GB free space minimum.
\item Network: Stable internet connectivity for cloud services.
\end{itemize}

\subsection{Software Requirements}
\begin{itemize}
\item Operating System: Windows/Linux/macOS.
\item Runtime: Node.js (LTS), npm.
\item Frameworks: Next.js, React, TypeScript.
\item Database and Auth: Supabase (PostgreSQL + Auth + Realtime).
\item Testing: Jest and Playwright.
\item Deployment: Vercel (or equivalent cloud platform).
\end{itemize}

\subsection{Feasibility Study}
\begin{enumerate}
\item Technical feasibility: High, using mature open-source tools.
\item Economic feasibility: Favorable due to free-tier cloud and open-source stack.
\item Operational feasibility: Strong fit for institutional processes.
\item Schedule feasibility: Achievable with phased implementation.
\end{enumerate}

\section{System Analysis}

\subsection{Introduction}
This chapter analyzes end-to-end QGX-D1 behavior, actors, interactions, and process-level flow patterns.

\subsection{System Overview}
QGX-D1 is a role-based cloud application with modular service boundaries and secure request handling.

\subsubsection{User Registration and Authentication}
Registration and login are managed through validated forms and secure session handling.

\subsubsection{Role-Based Dashboard and Navigation}
Each role receives a dedicated dashboard with authorized modules and navigation routes.

\subsubsection{AI Query Generation and Processing}
AI requests are validated, processed through backend APIs, and returned in structured format.

\subsubsection{Batch User Creation Workflow}
Admin can create users in bulk with row-level validation and per-record status reporting.

\subsubsection{Activity Logs and Admin Monitoring}
Critical operations are recorded to support auditing and monitoring.

\subsubsection{Password Recovery and Account Security}
Token-based password reset with expiry and single-use rules is implemented.

\subsubsection{PWA and Offline Support}
The system supports installable behavior and offline fallback page rendering.

\subsubsection{Settings and Profile Management}
Users can maintain profile data and account preferences with validation.

\subsubsection{Error Handling and Reliability}
Error boundaries, retry strategies, and fallback responses improve reliability.

\subsection{Actors of the System}
Actors include Admin, Teacher, Student, Parent, and external services (AI provider, database/auth provider).

\subsection{Use Case Diagram and Description}

\subsubsection{Actors}
\begin{itemize}
\item Admin: governance, user management, monitoring.
\item Teacher: assessments, content support, class operations.
\item Student: participation, assessment attempts, progress tracking.
\item Parent: student monitoring and communication.
\end{itemize}

\subsubsection{Algorithm and Flowcharts}
Core workflow uses validation-first processing, role guard checks, transactional persistence, and response normalization.

\subsubsection{Use Case Diagram}
\placeholderfigure{use-case}{Figure 3.1: Use Case Diagram}

\subsection{Activity Diagram}
\placeholderfigure{activity}{Figure 3.2: Activity Diagram for Major Workflow}

\subsection{Sequence Diagram}
\placeholderfigure{sequence}{Figure 3.3: Sequence Diagram for Authentication and Access}

\subsection{Class Diagram}
\placeholderfigure{class-diagram}{Figure 3.4: Class Diagram (Logical Model)}

\subsection{Component Diagram}
\placeholderfigure{component}{Figure 3.5: Component Diagram of QGX-D1}

\section{System Design}

\subsection{Introduction}
System design defines architecture, modules, data flow, database structure, security policy, and user experience decisions.

\subsection{System Architecture}
QGX-D1 follows layered architecture with presentation, application, and data/service layers.

\subsection{System Modules}
Modules include Authentication, Authorization, Dashboard, AI API, Batch User Management, Notifications, PWA, and Recovery.

\subsection{Data Flow Description}
Data flows from UI to API validation, role authorization, database persistence, and response delivery.

\subsection{Database Design}
The schema includes users, roles, sessions, tests, questions, attempts, notifications, and activity logs with relational integrity.

\subsection{Security Design}
Security controls include secure sessions, role validation at all layers, schema-based input checks, and immutable audit logs.

\subsection{AI-Assisted Workflow and User Experience Design}
AI support is constrained by governance checks, timeout handling, and user-visible result control.

\begin{table}[H]
\centering
\caption{Table 4.1: System Design Summary}
\label{tab:design-summary}
\begin{tabularx}{0.92\textwidth}{|l|X|}
\hline
\textbf{Design Area} & \textbf{Implementation Direction} \\
\hline
Architecture & Layered, role-centric, cloud-native \\
\hline
Security & Multi-layer checks with auditability \\
\hline
Data Model & Normalized PostgreSQL with RLS \\
\hline
AI Integration & Controlled API pipeline with validation \\
\hline
UX Direction & Role-specific dashboards and guided workflows \\
\hline
\end{tabularx}
\end{table}

\section{Implementation}

\subsection{Introduction}
This chapter presents practical implementation steps, technology choices, module realization, and deployment execution.

\subsection{Technologies Used}

\subsubsection{Backend Technologies}
Next.js API routes, TypeScript, Supabase services, validation libraries, and middleware controls.

\subsubsection{Frontend Technologies}
React components, Next.js app router, responsive CSS, and role-aware client logic.

\subsubsection{AI and NLP Technologies}
GROQ-based language model APIs are integrated through controlled backend handlers.

\subsubsection{Testing and QA Tooling}
Jest, Playwright, and static analysis tools are used to validate behavior and quality.

\subsection{System Implementation Overview}
Implementation follows modular coding, incremental integration, and continuous validation.

\subsection{Module Implementation}

\subsubsection{Authentication Module}
Implements registration, login, secure sessions, and reset tokens.

\subsubsection{Authorization and Role Management Module}
Provides role guards, permission mapping, and access control checks.

\subsubsection{AI API Module}
Implements input sanitization, model invocation, and response normalization.

\subsubsection{Batch User Management Module}
Supports bulk onboarding with validation and partial-failure reporting.

\subsubsection{Dashboard and Navigation Module}
Builds role-specific menu structures and route-level authorization.

\subsubsection{Notification and Toast Module}
Implements real-time and in-app feedback messaging.

\subsubsection{PWA and Service Worker Module}
Adds manifest, service worker registration, and offline page handling.

\subsubsection{Error and Recovery Module}
Adds error boundaries, retry logic, and fallback responses.

\subsection{Core Processing Algorithm}

\subsubsection{API Layer and Request Pipeline}
Requests pass through authentication, role checks, schema validation, business logic, and response serialization.

\subsubsection{Data Validation and Business Logic}
Input contracts are enforced before persistence; invalid input is rejected deterministically.

\subsubsection{Frontend Orchestration}
Client state and API orchestration are organized by feature modules.

\subsubsection{Database Schema and Persistence}
Relational schema and policy checks support secure and consistent persistence.

\subsection{Project File Structure}
The project structure contains app routes, reusable components, API modules, tests, scripts, and deployment configuration files.

\subsection{Installation and Execution}
\begin{enumerate}
\item Install dependencies using npm.
\item Configure environment variables for Supabase and AI APIs.
\item Run local development server.
\item Execute test suites (unit, integration, and E2E).
\item Build and deploy to cloud environment.
\end{enumerate}

\subsection{Implementation Challenges}
Challenges included role consistency, asynchronous API failure handling, and batch operation reporting. These were mitigated with centralized guards, timeout/retry policies, and record-level validation.

\section{Testing}

\subsection{Introduction}
Testing verifies correctness, security, usability, and reliability across all key workflows.

\subsection{Testing Methodology}
The strategy combines unit, integration, and end-to-end testing with security and accessibility validation.

\subsection{Unit Testing}
Unit tests cover utilities, validation logic, and role permission behavior.

\subsection{Integration Testing}
Integration tests validate API plus database behavior and role-restricted operations.

\subsection{End-to-End Testing}
E2E tests validate complete journeys including registration, login, dashboard access, and admin operations.

\subsection{Security and Accessibility Testing}
Security tests include unauthorized access rejection and malformed payload checks. Accessibility checks include semantic structure and interaction clarity.

\subsection{Test Results Summary}
Overall testing indicates strong pass rates for core workflows with known edge cases tracked for iterative improvement.

\section{Result Analysis and Output Screens}

\subsection{Introduction}
This chapter summarizes observed outputs, interface evidence, and system-level performance outcomes.

\subsection{Landing and Login Interface}
The landing and login interfaces provide clear identity entry paths and secure authentication behavior.

\subsection{Registration and Password Recovery Screens}
Registration and reset flows support validated onboarding and account recovery controls.

\subsection{Dashboard and Navigation Screens}
Role-specific dashboards provide task-focused navigation and controlled module visibility.

\subsection{AI Query and Response Screens}
AI query screens provide guided input and structured output with user-visible controls.

\subsection{Batch User Creation and Admin Screens}
Admin screens support bulk import and status reporting for onboarding operations.

\subsection{API and Performance Results}
API operations demonstrate acceptable latency under expected academic workflow load.

\subsection{Security and Reliability Results}
Role protection, request validation, and error handling show stable and secure behavior.

\subsection{Output Screens}
\placeholderfigure{output1}{Figure 7.1: Landing Interface Output}
\placeholderfigure{output2}{Figure 7.2: Role Dashboard Output}
\placeholderfigure{output3}{Figure 7.3: AI Query Output}

\section*{Conclusion}
QGX-D1 demonstrates that a modern, role-based, cloud-native LMS can deliver secure and maintainable academic operations with practical deployment readiness.

\section*{Future Scope}
Future scope includes adaptive learning analytics, multi-tenant architecture, deeper offline synchronization, and expanded AI-assisted personalization.

\newpage
\section*{Bibliography}
\small
\begin{thebibliography}{99}

\bibitem{anderson2022} Anderson, H., and Lewis, G., ``Microservices vs. Monolithic Architecture in Educational Technology,'' \textit{ACM Computing Surveys}, vol. 54, no. 8, pp. 1--35, 2022.

\bibitem{davis2023} Davis, P., Martinez, R., and Thompson, L., ``Integration of Large Language Models in Educational Technology: Benefits, Risks, and Governance Strategies,'' \textit{Journal of Educational Computing Research}, vol. 61, no. 4, pp. 856--879, 2023.

\bibitem{kumar2022} Kumar, S., Sharma, M., and Verma, R., ``Progressive Web Applications in Educational Contexts: Accessibility and Usability,'' \textit{Journal of Web Engineering}, vol. 21, no. 1, pp. 78--95, 2022.

\bibitem{lee2023} Lee, M., Chen, Y., and Wong, S., ``Role-Based Access Control in Educational Information Systems: A Systematic Review,'' \textit{Computers \& Education}, vol. 189, pp. 104--116, 2023.

\bibitem{nextjs} Next.js Documentation, Vercel. [Online]. Available: \url{https://nextjs.org/docs}

\bibitem{playwright} Playwright Documentation, Microsoft. [Online]. Available: \url{https://playwright.dev/docs}

\bibitem{supabase} Supabase Documentation, Supabase. [Online]. Available: \url{https://supabase.com/docs}

\bibitem{typescript} TypeScript Handbook, Microsoft. [Online]. Available: \url{https://www.typescriptlang.org/docs}

\end{thebibliography}

\newpage
\section*{Annexure}
\subsection*{Annexure A: Base Paper}
Attach the selected base research paper used for problem definition and system direction.

\subsection*{Annexure B: Published Paper}
Attach the published or submitted project paper based on QGX-D1.

\subsection*{Annexure C: Certificates}
Attach project completion certificate, internship/training certificate (if any), and publication certificate.

\end{document}
